\section{Discussion and Conclusion}
\label{sec:discussion}

Event-spline policies benefit more from clause-aligned training than waypoint policies under an identical clause schedule. This interaction repeats across three seeds. The complete interface factorial shows the resulting trade-off: adapting to scheduled clauses improves their execution but reduces performance with a static whole-task caption. Order interventions show a corresponding change in the first completed subgoal, while RoboCasa extends the label comparison to longer manipulation sequences.

Several boundaries remain. The head comparison uses one VLA backbone and two LIBERO tasks, and combines event alignment, compactness, and duration prediction. Reversed-order completion remains difficult, so the current results do not establish general recombination of learned skills. Fixed clause clocks also assume predictable phase timing; the action-release condition relaxes that assumption only for tasks with a suitable gripper event.

Separating a predicted path from its duration lets the executor retime a motion without retraining the policy. CALVIN shows improved completion and pace for each frozen spline policy. An equally retimed waypoint is a strong control, establishing that generic temporal resampling is shared by both representations. The learned event duration provides the additional interface: it selects motions for compression and predicts when a new observation is useful. Together, the studies evaluate one action representation for selecting subgoals and controlling the timing of their motion.
